\documentclass[11pt]{article}
\usepackage[margin=1in]{geometry}
\usepackage{amsmath,amssymb,amsthm}
\usepackage{hyperref}
\hypersetup{colorlinks=true,urlcolor=blue,linkcolor=blue}
\newtheorem{theorem}{Theorem}
\title{A fourth configuration with stationary orthocenter\\
\large MathDB 365004 (duplicate formulation 372155)}
\author{Principia Math In Progress}
\date{21 September 2026}
\begin{document}
\maketitle
\section*{Task 1: a full counterexample}
The classification conjecture of Garcia, Helman and Reznik
\cite[Conjecture 1, p.6]{special} asserts that a Poncelet triangle family
can have a stationary orthocenter only in the focal-$X_4$, MacBeath, or
Dual configurations. The following counterexample has nested,
noncircular ellipses and consists entirely of acute triangles.

\begin{theorem}
The two ellipses
\[
 E:\ \frac{x^2}{4}+y^2=1,\qquad
 K:\ \frac{(x-31/20)^2}{4/25}+\frac{y^2}{63/1600}=1
\]
support a full Poncelet family of triangles. Every triangle has
orthocenter $H=(31/16,0)$. This conic pair is none of the three
configurations in the conjecture.
\end{theorem}

\noindent\textbf{Credit and prior work.}
The general orthocenter-locus semiaxis formula is already
\cite[Theorem 3]{general}; the example can be obtained by specializing
that formula. We explicitly credit it and give a direct algebraic
verification below. The campaign assigns \textbf{zero new credit}
pending inspection of the identified journal final of the original
source \cite{original}. No new general locus formula is claimed.

\section*{1. Nesting and the full family}
Apply $(x,y)=(2X,Y)$ to the unit circle. The inverse image of $K$ is
\[
 K_0:\quad (X,Y)=
 \left(\frac{31}{40}+\frac15\cos t,\frac{3\sqrt7}{40}\sin t\right).
\]
With $v=\cos t$ its squared distance to the origin is
\[
 X^2+Y^2=\frac{v^2+496v+1024}{1600}
       \leq\frac{1521}{1600}<1 .
\]
The polynomial is increasing on $[-1,1]$. Thus the caustic is strictly
nested, and both its semiaxes are positive.

\newpage
Put $f=4/5$, $g=3/4$, $s=f+g=31/20$ and $p=fg=3/5$. For every
$|\lambda|=1$, let $a,b,c$ be the roots of
\begin{equation}\label{cubic}
 P_\lambda(z)=z^3-(s+p\lambda)z^2+(p+s\lambda)z-\lambda.
\end{equation}
The equation $P_\lambda(z)=0$ is equivalent on the unit circle to
\[
 B(z)=\lambda,\qquad
 B(z)=\frac{z(z-f)(z-g)}{(1-fz)(1-gz)}.
\]
This degree-three finite Blaschke product has strictly positive
argument derivative
$1+(1-f^2)/|e^{it}-f|^2+(1-g^2)/|e^{it}-g|^2$.
Its total argument change is $6\pi$. Consequently every $\lambda$
has exactly three distinct roots on the unit circle. The triangles
with vertices $Z(a),Z(b),Z(c)$, where
$Z(z)=(3z+z^{-1})/2$, are nondegenerate and inscribed in $E$.

\section*{2. Every side touches the same caustic}
For a pair $a,b$ write $q=a+b$, $r=ab$.
Vieta's formulas from \eqref{cubic} imply
\[
 c(1-pr)=s-q,\qquad r-p=c(sr-q).
\]
Here $1-pr\ne0$, since $|pr|=3/5$. Eliminating $c$ gives
\begin{equation}\label{pair}
 q^2-\frac{31}{20}(1+r)q+\frac35(1+r^2)
                         +\frac{417}{400}r=0.
\end{equation}
The chord of the unit circle through $a,b$ has equation
$z+r\overline z=q$. Write $r=e^{2i\theta}$ and
$q=2d e^{i\theta}$ with $d\in\mathbb R$. This is the real line
$X\cos\theta+Y\sin\theta=d$.
Its tangency condition to $K_0$ is
\[
 \left(d-\frac{31}{40}\cos\theta\right)^2
       =\frac1{25}\cos^2\theta+\frac{63}{1600}\sin^2\theta.
\]
Multiplication by $4r\cdot1600$ converts this condition into
\[
 1600q^2-2480q(r+1)+897(r+1)^2+63(r-1)^2=0.
\]
This is precisely 1600 times \eqref{pair}.
Every side therefore touches $K_0$, and its affine image touches $K$.

The caustic is inside each triangle. At $\lambda=1$, the roots are
$1$ and $23/40\pm3i\sqrt{119}/40$, giving the physical triangle
\[
 (2,0),\quad
 \left(\frac{23}{20},\frac{3\sqrt{119}}{40}\right),\quad
 \left(\frac{23}{20},-\frac{3\sqrt{119}}{40}\right).
\]
Its interior contains the caustic center $(31/20,0)$. As $\lambda$
varies, roots remain distinct, and a tangent line cannot cross the
center of a nondegenerate ellipse. Continuity preserves the interior
side for all three supporting lines. Thus all these triangles
circumscribe $K$. They form a full Poncelet family: every point of $E$
occurs as a vertex, by choosing $\lambda=B(z)$ for its preimage $z$.

\newpage
\section*{3. All altitudes pass through the fixed point}
Put $A=3/2$, $D=1/2$, $h=31/16$, $L=abc$, $s_1=a+b+c$ and
$s_2=ab+bc+ca$. For unit $a,b,c$,
$Z(a)=Aa+D/a$ and $\overline{Z(a)}=A/a+Da$.
Direct expansion gives
\begin{align*}
 &(h-Z(a))(\overline{Z(b)}-\overline{Z(c)})
 +(h-\overline{Z(a)})(Z(b)-Z(c))\\
 &\quad=\frac{b-c}{L}
 \left[(DL-Aa)h+(AL-Da)h-2ADLa
       +(A^2+D^2)s_1a-(A^2+D^2)s_2+2AD\right].
\end{align*}
Substitute $L=\lambda$, $s_1=s+p\lambda$,
$s_2=p+s\lambda$. The bracket is identically zero:
its $a$, $\lambda a$, $\lambda$, and constant coefficients all vanish.
The left side is twice the real Euclidean dot product of
$H-Z(a)$ with $Z(b)-Z(c)$. Thus $H$ lies on the altitude from $Z(a)$;
cyclically it lies on all three. It is their orthocenter.

\section*{4. Exclusion of the three configurations}
The definitions in \cite[Sections 3--4]{special} distinguish the three
families by their actual Euclidean conic configurations:

\begin{center}
\begin{tabular}{ll}
Configuration & Reason our pair is different\\ \hline
MacBeath & The outer conic is a circle; $E$ has semiaxes $2,1$.\\
Focal-$X_4$ & The inner conic is a circle; $4/25\ne63/1600$.\\
Dual & The two ellipses are concentric; their centers here differ.
\end{tabular}
\end{center}

These distinctions are preserved by similarities. In addition,
$H$ is not a focus of $E$, since $(31/16)^2\ne3$.
Finally,
\[
             \frac{(31/16-31/20)^2}{4/25}
                    =\frac{961}{1024}<1.
\]
So $H$ lies strictly inside $K$ and therefore inside every family
triangle. Every triangle is acute. Neither degeneracy nor an obtuse
exception is involved. This completes the counterexample.

\section*{Prior formula and audit status}
For outer semiaxes $a_0,b_0$, the known formula
\cite[Theorem 3]{general} has the semiaxis numerator
\[
 \sigma=\left|fg(a_0^2+b_0^2)-(a_0^2-b_0^2)\right|.
\]
Our choice gives $\sigma=|(3/5)5-3|=0$. Thus the example is consistent
with that earlier theorem. Its application to the later proposed
three-family classification is the subject of this note.

\newpage
The full current 24-page original preprint, the complete 7-page
classification preprint, and the complete 18-page general-locus preprint
were read. The 32-page final journal article corresponding to the
original preprint remains unread; the publisher provides a subscription
preview. Source and priority eligibility therefore remain pending.
MathDB 365004 and 372155 are the same problem, not two credits.

The accompanying verifier checks nesting, the exact cubic and tangency
identities, and the universal altitude identity. It also checks
48 configurations independently in Euclidean coordinates at 90 decimal
digits. These numerical checks supplement the proof.
Review is by the authoring agent only.

\section*{Proposed follow-up Tasks 2--5}
\begin{enumerate}
\item \textbf{Task 2.} Find a synthetic explanation for the stationary
orthocenter of this pair.
\item \textbf{Task 3.} Classify all nested ellipse pairs with stationary
orthocenter, up to similarities, using the credited general formula.
\item \textbf{Task 4.} Describe the exact region of attainable stationary
orthocenters for a prescribed outer ellipse.
\item \textbf{Task 5.} Determine the polar-circle invariants of every
family in that classification.
\end{enumerate}

\begin{thebibliography}{9}
\bibitem{special} R. A. Garcia, M. Helman, D. Reznik,
\emph{Four special Poncelet triangle families about the incircle}.
\href{https://arxiv.org/abs/2512.17440}{arXiv:2512.17440v2} (2025).
\bibitem{general} R. A. Garcia, M. Helman, D. Reznik,
\emph{Poncelet triangles: conic loci of the orthocenter and of the
isogonal conjugate of a fixed point}.
\href{https://arxiv.org/abs/2508.02368}{arXiv:2508.02368v3} (2025).
\bibitem{original} R. A. Garcia, M. Helman, D. Reznik,
\emph{Harmonious loci of Poncelet triangles about the incircle and their
degeneracies}. \href{https://arxiv.org/abs/2409.19464}{arXiv:2409.19464v7}
(2025). Journal final under the title \emph{Graceful loci},
Beitr.\ Algebra Geom.\ \textbf{67} (2026), 667--698,
\href{https://doi.org/10.1007/s13366-026-00825-1}{doi:10.1007/s13366-026-00825-1};
final unread.
\end{thebibliography}
\end{document}
